\documentclass[11pt,a4paper]{article}
\usepackage[T1]{fontenc}
\usepackage{lmodern,amsmath,amssymb,amsthm,booktabs,microtype}
\usepackage[margin=27mm]{geometry}
\usepackage[colorlinks=true,linkcolor=blue!50!black,citecolor=blue!50!black,urlcolor=blue!50!black]{hyperref}
\usepackage{xcolor}
\setlength{\emergencystretch}{2em}
\newtheorem{theorem}{Theorem}[section]
\newtheorem{proposition}[theorem]{Proposition}
\newtheorem{lemma}[theorem]{Lemma}
\newtheorem{corollary}[theorem]{Corollary}
\theoremstyle{remark}\newtheorem{remark}[theorem]{Remark}
\newcommand{\Q}{\mathbb Q}\newcommand{\Z}{\mathbb Z}
\title{Arithmetic dependencies of the construction with 29 variables}
\author{}\date{8 November 2026}
\begin{document}\maketitle
\begin{abstract}
We identify the arithmetic results used by the finite-algebra construction
with 29 total variables. Its four stated arithmetic inputs follow
from the specified elliptic-group, height and pole theorems in the source
manuscript and Sun's simultaneous prime-pattern theorem. In particular,
the required avoidance at seven, compatibility with an arbitrary fixed
congruence progression, and realization with zero preparation parameters
require no stronger global existence assertion. The deductions are given
explicitly. The deep arithmetic theorems themselves remain cited inputs.
\end{abstract}

\section{The statements being used}

Fix
\[
 K=\Q(\sqrt2),\quad D=11-6\sqrt2,\quad E:y^2=x^3-Dx,
 \qquad
 P=\left(\frac{57-28\sqrt2}{8},\frac{-660+417\sqrt2}{32}\right).
\]
Heights are absolute logarithmic Weil heights, and rational valuations
are normalized by $v_p(p)=1$. We count the distinguished rational input
as one of the variables of a polynomial.

The source manuscript used here is Version~9 of \cite{GlazerV9}, at commit
\begin{center}\small\texttt{14e8f601fe0d9173d4a5f084aedc3e07066a7688}.\end{center}
The relevant file is \texttt{manuscripts/quartic\_v9.tex}.
Sun's theorem is Theorem~4.1 of the explicitly versioned preprint
\cite{Sun2026}, arXiv:2610.20691v2. File hashes and theorem identifiers
are recorded in the accompanying source certificate. These pins fix
the statements whose implications are proved below. For the underlying
restricted-converse argument we incorporate the Version~11 correction
\cite{GlazerV11}: the Kato classes are defined on the full index set
needed by the 1999 theorem, and their norms are identified with the
retained cyclotomic class. This correction changes the proof supporting
an inherited input, not the statement of the input used here.

\begin{enumerate}
\item The elliptic source gives $2E(K)=\Z P$ and
      $h(x([n]P))\ge n^2-6$ for $n\ne0$.
      The relevant source identifiers are
      \nolinkurl{subsec:v5-point-source} and
      \nolinkurl{eq:o8-indexheight}.
\item Put
\[
 B_* = \prod_{p\le76176}p,
 \qquad H_{19}=2^{110}210^{1800}B_*^{19}+128.
\]
      For $t=u/v$ in lowest terms, $v>0$, and
      $t\notin\{0,1,2401,-343/32\}$, let $M_5(t)$ be the product
      of primes $p>76176$ for which at least one of
\[
 v_p(v),\quad v_p(u),\quad v_p(u-v),\quad
 v_p(u-2401v),\quad v_p(32u+343v)
\]
      is positive and odd. The source theorem
      \nolinkurl{thm:quartic-height19} gives
      $h(t)\le H_{19}M_5(t)^{19}$.
\item The pole results are the source obstruction, simultaneous
      realization and zero-preparation lemmas, respectively:
      \begin{quote}\small
      \nolinkurl{hyp:o8-pole},\\
      \nolinkurl{lem:original-seven-simultaneous-realization},\\
      \nolinkurl{v5:preparation-zero-bridge}.
      \end{quote}
      The last of these has exactly the zero-preparation conclusion
      used in Section~\ref{sec:poles}.
\item The prime-pattern result is Sun's Theorem~4.1, whose implications
      for the prescribed congruences and the elliptic reduction orders
      are proved in Section~\ref{sec:avoidance}. The same quantified
      statement appears as \nolinkurl{thm:v5-simultaneous-primes}
      in the source manuscript.
\end{enumerate}

The first two entries involve the identical normalized point and the
identical height constant. The exact integer calculation
\[
                      2^{2093447}<H_{19}<2^{2093448}
\]
is a sharper evaluation of this fixed constant; it does not strengthen
the source height inequality. Likewise the implication with exponent
24 follows from the exponent-19 result because $M_5(t)\ge1$.

\section{Finite avoidance and the common congruence progression}
\label{sec:avoidance}

For a Gaussian integer $\pi$, set
\[
 d_-(\pi)=N_{\Q(i)/\Q}(\pi-1),\qquad
 d_+(\pi)=N_{\Q(i)/\Q}(\pi+1).
\]
Sun's theorem permits any fixed finite list of distinct integral
coefficients, any fixed divisibility of the square-denominator parameter,
and any fixed lower bounds for the prime values and for the slope.
Its avoidance modulus is odd, squarefree, and prime to seven.
The fixed Gaussian residues are $9$ modulo $16$ and $1+3i$ modulo seven;
at every prime dividing the avoidance modulus the residue is $i$.
The theorem excludes every odd common divisor between the two order
norms at the same position and between any two order norms at distinct
positions.

\begin{lemma}\label{lem:arbitrary-avoidance}
Let $L\ge1$ be divisible by $64$, and fix distinct integers
$c_1,\ldots,c_k$, with $k\ge2$. Prescribe $q\ge1$ and real numbers
$Y,H_0>0$. Sun's Theorem~4.1 supplies positive integers $u,w$ such that
\[
 q\mid w,\qquad u/w^2>H_0,\qquad p_j=u+c_jw^2>Y
\]
are primes, with Gaussian generators whose order norms satisfy
\[
 \begin{gathered}
 v_2(d_-(\pi_j))=6,\qquad v_2(d_+(\pi_j))=2,\\
 \gcd(d_-(\pi_j),d_+(\pi_h))=4\qquad(1\le j,h\le k),
 \end{gathered}
\]
and are all prime to the odd part of $L$. The primes may simultaneously
be required to satisfy $p_j\equiv1\pmod{8\cdot3\cdot17}$.
\end{lemma}
\begin{proof}
Take the avoidance modulus to be the product of the odd primes dividing
$3\cdot17\cdot L$, with seven omitted. It is squarefree and prime
to seven, as required by Sun's theorem. At each of its prime factors,
the two adjacent norms of $i$ are both two. They are therefore units.
At seven the fixed residue $1+3i$ instead gives
\[
       d_-(1+3i)=9\equiv2\pmod7,\qquad
       d_+(1+3i)=13\equiv6\pmod7.
\]
Thus seven is avoided even when it divides $L$.

Write $\pi=9+16a+16bi$. Direct calculation gives
\[
 d_+(\pi)\equiv4\pmod8,
 \qquad d_-(\pi)\equiv64\pmod{128}.
\]
The exact powers of two in the statement follow. The exclusion of odd
common factors supplied by Sun's theorem then gives every stated
cross-sign gcd, including a position paired with itself.
Finally the Gaussian norm is one modulo eight, and its residues at
three and seventeen are the norm of $i$, namely one. These are precisely
the required congruences on $p_j$.
\end{proof}

\begin{corollary}\label{cor:compatible-progressions}
In the situation of Lemma~\ref{lem:arbitrary-avoidance}, put
\[
 U=\operatorname{lcm}_j\bigl(d_-(\pi_j)/4\bigr),\qquad
 V=\operatorname{lcm}_j\bigl(d_+(\pi_j)/4\bigr).
\]
Then
\[
              \gcd(U,V)=1,\qquad\gcd(U,L)=16,
              \qquad\gcd(V,L)=1.
\]
For every integer $h_*$ divisible by $64$, the congruences
\[
                 h\equiv0\pmod U,\qquad
                 h\equiv1\pmod V,\qquad
                 h\equiv h_*\pmod L
\]
have an arithmetic progression of solutions.
\end{corollary}
\begin{proof}
Division by four makes every minus order have two-adic valuation four
and every plus order odd. The cross-sign gcds make their least common
multiples coprime. Lemma~\ref{lem:arbitrary-avoidance} excludes all odd
factors of $L$ from both least common multiples. The displayed gcds
follow. The only nontrivial pairwise compatibility condition is
$h_*\equiv0\pmod{16}$, which holds. The generalized Chinese remainder
theorem now supplies the progression.
\end{proof}

The reduction orders at the two places of $K$ above $p_j$ are the
unordered pair $d_-(\pi_j),d_+(\pi_j)$, by the source lemma
\nolinkurl{lem:v5-two-reductions}. Both reductions have full rational
two-torsion. For a finite group of elliptic-curve points of order $d$,
full two-torsion implies that the subgroup of doubles has order $d/4$.
Every double is therefore killed by $d/4$. This justifies the divisions
by four in Corollary~\ref{cor:compatible-progressions}; the source point
$P$ is itself a double.

\begin{lemma}\label{lem:rational-contacts}
The preceding statements also apply to any fixed finite list of
rational contact coefficients, with a common square denominator for
the slope and with arbitrarily deep fixed denominator divisibility.
\end{lemma}
\begin{proof}
For coefficients $\gamma_j\in\Q$, choose a positive integer $D_0$
so that $D_0^2\gamma_j$ are integers, and remove repeated coefficients.
Apply Lemma~\ref{lem:arbitrary-avoidance} with
$c_j=D_0^2\gamma_j$ and put $v=D_0^2w^2$.
Then $u+c_jw^2=u+\gamma_jv$ and $u/v>D_0^{-2}H_0$.
Enlarge the prescribed lower bound before applying the lemma to obtain
any desired slope bound. By choosing $q$ divisible by suitable powers
of a fixed finite set of primes, one obtains any prescribed lower
bounds on their valuations in $v$. Enlarge $Y$ to exclude any fixed
finite list of primes. None of these choices changes the quantified
form of Sun's theorem.
\end{proof}

\begin{lemma}\label{lem:real-progression}
Let $G$ be the closure of a cyclic subgroup generated by a point $Q$
in a compact real Lie group, and let $G^0$ be its identity component.
Suppose $L$ is divisible by the order of $G/G^0$. Any arithmetic
progression $h\equiv h_*\pmod L$, with an additional compatible
finite collection of congruences, is dense after multiplication by
$Q$ in the component $h_*Q+G^0$. Arbitrarily large positive indices
can be used.
\end{lemma}
\begin{proof}
The compatible congruences give one coset modulo an integer $L'$
divisible by $L$. A compact connected abelian Lie group is a torus,
and multiplication by any nonzero integer is onto it. It follows that
$[L']G=G^0$. Thus the two-sided progression has the stated closure.
In a compact group, the closure of a positive cyclic orbit is also
a group: positive return times approach the identity and consequently
approximate the negative multiples. Removing a finite initial segment
does not change the closure. This proves the positive-index assertion.
\end{proof}

For the construction, apply this lemma to the closure of
$(P,P^\sigma)$ in the product of the two real elliptic groups.
The scalar point, the four integer root labels, the real compact
neighborhoods, all contact coefficients, and the modulus $L$ are fixed
before the prime-pattern theorem is used. Corollary~\ref{cor:compatible-progressions}
then supplies the final shared index without disturbing those
neighborhoods. No estimate uniform in the coefficients, $q$, $L$, or
the number of contact positions is required.

\section{Pole realization with zero preparation}
\label{sec:poles}

For $b\in\Q^\times$, put $B=17b$. The zero-preparation predicate is
\[
 \begin{aligned}
 p^2&=B(1-e)(1+3e),\\
 h_0&=e(B-b^2),\\
 Bq^2&=(s^2-Bh_0)(s^2+3Bh_0),
 \end{aligned}
\]
with rational $e,p,q,s$; $h_0$ is the displayed expression.
No nonzero condition is imposed on these four witnesses.

The source lemma \nolinkurl{v5:preparation-zero-bridge} states exactly
the following two implications. A solution excludes negative odd
valuations of $b$ outside $\{2,3,17\}$. Conversely a solution exists if
the reduced denominator of $b$ is a square and
\[
 \begin{gathered}
 0<b<17,\qquad b\in\Q_2^{\times2}\cap\Q_3^{\times2}
                              \cap\Q_{17}^{\times2},\\
 v_2(b)\ge3,\qquad v_3(b)\ge1,\qquad v_{17}(b)\ge2.
 \end{gathered}
\]
These are precisely the hypotheses verified for each of the four
normalized pole inputs in the finite-algebra construction.

The special role of zero preparation can also be checked directly in
the source argument. With preparation parameter zero, the auxiliary
quantity there is
\[
                    B(B-b^2)=B^2(1-b/17).
\]
It is positive and is a square in the three specified completions,
by the displayed bounds. If $\mu,\theta$ are the positive squarefree
representatives of $B$ and this quantity, their supports are disjoint,
$\mu\equiv17\pmod{24}$, $\theta\equiv1\pmod{24}$, and each is a
local square at primes dividing the other. To see the last condition,
an odd prime in the support of $\mu$ has positive odd valuation on $b$,
so $1-b/17$ is a square unit; a prime in the support of $\theta$ outside
$2,3,17$ must have $v_p(b)=0$ and $b\equiv17\pmod p$, making
$B\equiv17^2\pmod p$. The conditions at $2,3,17$ were already imposed.
The source simultaneous-realization lemma therefore applies without
an additional hypothesis. Its deep restricted-converse and rank
dependencies are retained in this use.

\begin{lemma}\label{lem:conic-complete}
For $B\in\Q^\times$, the first conic is covered by the two formulas
\[
 e=\frac{\tau^2+(4\varepsilon-1)B}{\tau^2+3B},\qquad
 p=\frac{-4B\tau(1-\varepsilon)}{\tau^2+3B},\qquad
 \varepsilon\in\{0,1\},\quad\tau^2+3B\ne0.
\]
Thus its parametrization imposes no additional restriction on the
pole-realization theorem.
\end{lemma}
\begin{proof}
Substitution verifies the conic equation. If $e\ne1$, take
$\varepsilon=0$ and $\tau=p/(e-1)$. The conic equation gives
$\tau^2+3B=4B/(1-e)\ne0$ and recovers the displayed point.
If $e=1$, then $p=0$, and $\varepsilon=1,\tau=0$ works.
The case $e=-1/3,p=0$ is included in the first chart.
\end{proof}

The two structural coordinates attached to different pole inputs are
independent. Once all four inputs satisfy the source hypotheses, choose
each rational certificate separately and use Lemma~\ref{lem:conic-complete}.
Zero second-square roots are included by the explicit finite support
split in the finite-algebra construction. No common structural witness
for distinct pole inputs is asserted or needed.

\section{The resulting implication}

\begin{theorem}\label{thm:dependency-conclusion}
Assume the cited elliptic-group, index-height, contact-height and
pole-realization theorems of the pinned source manuscript, and
Sun's Theorem~4.1. Then the four arithmetic inputs of the construction
with 29 variables hold. Consequently its proved rational-source
argument and finite-algebra compilation give a polynomial $P$ such that
\[
 t\in\Z\quad\Longleftrightarrow\quad
 \exists\mathbf y\in\Q^{28}:P(t,\mathbf y)=0.
\]
No additional global arithmetic existence theorem is required.
\end{theorem}
\begin{proof}
The normalized group and index height are identical to those in the
source. The contact-height constant is the same integer, evaluated
exactly. Section~\ref{sec:poles} identifies the needed pole theorem and
checks its parametrization and independent simultaneous use.
Section~\ref{sec:avoidance} supplies the prescribed prime choices,
all fixed odd avoidance, the exact congruence compatibility and the
unchanged real closure. These are all four arithmetic inputs.

The remaining steps are proved within the finite-algebra construction:
the odd-product height identity, the two-parameter point charts,
the scalar input normalization, the elementary Lucas growth estimate,
the fixed-prime scalar size condition, and the one-coordinate
elimination of the finite roots and nonzero guards. Their soundness and
completeness therefore apply with precisely the arithmetic inputs just
verified.
\end{proof}

The fixed-prime scalar normalization only adds its prime and the finitely
many denominator primes to the fixed avoidance set. The preceding
argument permits precisely this enlargement before Sun's theorem is
applied. It therefore does not alter the arithmetic dependency list.

The theorem is an implication from explicitly identified mathematical
results. It does not replace proofs of those deep source results by
finite computation, and it does not claim an independent verification
of the restricted converse, contact-height theorem, or analytic
prime-pattern theorem. The computational certificates verify the
literal formulas, source bindings, and finite numerical deductions.

\begin{thebibliography}{9}\raggedright
\bibitem{GlazerV9} E.~Glazer and contributors,
\emph{Diophantine definition of the integers in the rationals},
Version~9, 7 November 2026, source commit
\texttt{14e8f601fe0d9173d4a5f084aedc3e07066a7688}.
\url{https://github.com/elliotglazer/diophantine-z-q/tree/14e8f601fe0d9173d4a5f084aedc3e07066a7688}.
\bibitem{GlazerV11} E.~Glazer and contributors,
\emph{Diophantine definition of the integers in the rationals},
Version~11, Kato index-set and norm-identification correction,
source commit \texttt{650b298a2a18fa956f8859f3ca80e9d7c8cedd82}.
\url{https://github.com/elliotglazer/diophantine-z-q/tree/650b298a2a18fa956f8859f3ca80e9d7c8cedd82/docs/manuscript-revisions/version11}.
\bibitem{Sun2026} Z.-W.~Sun,
\emph{86 unknowns for Hilbert's tenth problem over the rationals},
arXiv:2610.20691v2, 19 October 2026, especially Theorem~4.1.
\url{https://arxiv.org/abs/2610.20691v2}.
\end{thebibliography}
\end{document}
